\section{What determines the error: a decomposition and closed-form predictors}
\label{sec:theory}

Throughout, the randomness is the assignment $z$ only (design-based); potential outcomes are fixed. Let $\phi_i(z)\in\mathbb R^k$ be the feature map, $\hat\Sigma=\frac1n\sum_i\phi_i\phi_i^\top$ and $\hat\beta=\hat\Sigma^{-1}\frac1n\sum_i\phi_iy_i$. The estimator is $\hat\tau=\bar d^\top\hat\beta$ with $\bar d=\frac1n\sum_i[\phi_i(\mathbf 1)-\phi_i(\mathbf 0)]$ (for $\phi=(1,z,s)$: $\bar d=(0,1,1)$). Let $\Sigma=\mathbb E\hat\Sigma$, $\beta^\star=\Sigma^{-1}\frac1n\sum_i\mathbb E[\phi_iy_i]$ be the population least-squares coefficient under the design, $\varepsilon_i=y_i-\phi_i^\top\beta^\star$ and $\hat g=\frac1n\sum_i\phi_i\varepsilon_i$ (so $\mathbb E\hat g=0$).

\begin{proposition}[Decomposition]\label{prop:decomp}
Exactly, $\hat\tau-\tau=B+\bar d^\top\hat\Sigma^{-1}\hat g$ with the \emph{extrapolation bias} $B=\bar d^\top\beta^\star-\tau$, which does not depend on $n$ through any sampling effect. To first order ($\hat\Sigma\to\Sigma$) the second term has mean zero and variance at most $\kappa L/n$, where $L=\bar d^\top\Sigma^{-1}\bar d$ is the \emph{design leverage} of the all-treated-vs-all-control contrast and $\kappa=\lambda_{\max}\big(\Sigma^{-1/2}\,n\mathrm{Cov}(\hat g)\,\Sigma^{-1/2}\big)$ is the effective residual scale ($\kappa\approx\sigma^2$ for homoskedastic independent residuals). Hence $\mathrm{RMSE}^2\approx B^2+\kappa L/n$.
\end{proposition}

$L$ depends on the graph, design and feature map but not on the outcomes: it is the price of extrapolating from the exposures the design generates to the corners ``everyone treated'' and ``no one treated''. For $\phi=(1,z,s)$ (Chin-type estimator) it has a closed form.

\begin{proposition}[Leverage of the one-hop fraction model]\label{prop:lev}
Let $v=p(1-p)$, $w_{ic}=m_{ic}/d_i$ where $m_{ic}$ is the number of neighbours of $i$ in cluster $c$, $\rho=\frac1n\sum_i w_{i,c(i)}$ and $q=\frac1n\sum_i\sum_c w_{ic}^2$. Under independent cluster-level Bernoulli$(p)$ assignment (unit-level Bernoulli is the case of singleton clusters, $\rho=0$, $q=1/d_H$ with $d_H=(\frac1n\sum_i 1/d_i)^{-1}$ the harmonic-mean degree),
\[
L=\frac{1}{v}\Big[1+\frac{(1-\rho)^2}{q-\rho^2}\Big],\qquad\text{in particular}\qquad L_{\rm Bern}=\frac{1+d_H}{p(1-p)} .
\]
\end{proposition}

The interference penalty over the no-interference variance $1/v$ is $(1-\rho)^2/(q-\rho^2)$: $d_H$ for unit randomization, $\to0$ when all neighbours share the unit's cluster ($\rho\to1$; $z$ and $s$ collinear, as precise as DiM). So $\mathrm{RMSE}\propto\sqrt{(1+d_H)/(vn)}$ under unit randomization.

\begin{proposition}[Extrapolation bias under a nonlinear exposure response]\label{prop:extrap}
Let $Y_i(z)=a_i+\beta_iz_i+\gamma\,g(s_i)$ with $g(0)=0,\ g(1)=1$, and Bernoulli$(p)$ assignment (heterogeneous $\gamma_i$ with mean $\gamma$, independent of degree, give approximately the same $B$). For $\phi=(1,z,s)$,
\[
B=\gamma\,(\kappa_p-1),\qquad \kappa_p=\frac{\sum_i\mathrm{Cov}(g(s_i),s_i)}{\sum_i\mathrm{Var}(s_i)},\quad s_i\sim\mathrm{Bin}(d_i,p)/d_i ,
\]
and $\kappa_p\to g'(p)$ when all degrees grow, so that $B\to\gamma\,(g'(p)-1)$: \emph{the regression recovers the slope of the response at $p$, not the secant slope between $0$ and $1$.}
\end{proposition}

For $g(s)=s^k$ at $p=\tfrac12$, $g'(p)=k2^{1-k}$: $B=0$ for $k=2$ (tangent at $\frac12$ equals secant), $-\gamma/2$ for $k=4$, $\approx-0.29\gamma$ for $k=\frac12$; steep thresholds centred at $p$ give large positive bias. $B$ is not reduced by more data, nor (Sec.~\ref{sec:results}) by more flexible regressions when, under unit randomization with large degree, essentially no unit has exposure near $0$ or $1$.

\begin{proposition}[Heterogeneous effects: the estimand is variance-weighted]\label{prop:wt}
Let $Y_i(z)=a_i+\beta_iz_i+\gamma_is_i$ and Bernoulli assignment. For $\phi=(1,z,s)$, $\beta^\star$ yields
$\;\bar d^\top\beta^\star=\bar\beta+\sum_i\pi_i\gamma_i,\ \pi_i=d_i^{-1}/\sum_jd_j^{-1}$, whereas $\tau=\bar\beta+\bar\gamma$ (uniform weights). The bias is $B=\sum_i(\pi_i-\frac1n)\gamma_i$.
\end{proposition}

OLS weights unit effects by the variance $\propto1/d_i$ of the unit's exposure, so low-degree units dominate; if susceptibility correlates with degree the bias can be of the order of the effect itself, whatever $n$.

\begin{proposition}[Omitted multi-hop spillover]\label{prop:mh}
Let $Y(z)=(I-\lambda W)^{-1}(a+\beta z)$, $W=D^{-1}A$, $|\lambda|<1$, with Bernoulli assignment and $\phi=(1,z,s)$. Define $r_k=\frac1n\mathrm{tr}\,W^k$, $u_k=\frac1n\sum_{ij}(W^k)_{ij}W_{ij}$, $q=u_1$ and $\psi_k=r_k+u_k/q$. Then
\[ B=-\beta\sum_{k\ge2}\lambda^k(1-\psi_k)\;\approx\;-\beta\,\frac{\lambda^2}{1-\lambda}\quad(\psi_k\approx0).\]
\end{proposition}

$\psi_k$ is the share of $k$-step influence aliased with the included one-hop terms (return and closure probabilities); on tree-like graphs the regression misses all indirect effect beyond one hop. Under cluster randomization distant exposures correlate with $s$ and are partly absorbed (Sec.~\ref{sec:results}); we give no closed form for that case.

Proofs (Appendix~\ref{app:proofs}) are design-expectation statements; finite-$n$ and finite-cluster effects are quantified empirically.
